\documentclass[aspectratio=169,10pt]{beamer}
\usetheme{default}
\setbeamertemplate{navigation symbols}{}
\usepackage{graphicx}\usepackage[T1]{fontenc}\usepackage{amsmath,amssymb}
\usepackage{listings}
\definecolor{ncsured}{HTML}{CC0000}\definecolor{ncsudk}{HTML}{990000}
\definecolor{ncsugray}{HTML}{6D6D6D}\definecolor{ncsuteal}{HTML}{008473}
\definecolor{codebg}{HTML}{F5F5F5}
\setbeamercolor{structure}{fg=ncsured}
\setbeamercolor{frametitle}{fg=white,bg=ncsured}
\setbeamercolor{title}{fg=ncsudk}
\setbeamercolor{itemize item}{fg=ncsured}
\setbeamercolor{itemize subitem}{fg=ncsured}
\setbeamerfont{frametitle}{series=\bfseries}
\setbeamertemplate{frametitle}{\nointerlineskip\vskip2pt%
  \begin{beamercolorbox}[wd=\paperwidth,ht=2.4ex,dp=1ex,leftskip=.3cm]{frametitle}%
  \small\bfseries\insertframetitle\end{beamercolorbox}}
\setbeamertemplate{footline}{\hbox{\begin{beamercolorbox}[wd=\paperwidth,ht=2.2ex,dp=1ex,leftskip=.3cm,rightskip=.3cm]{}%
  \tiny Optimal Market Making for Cryptocurrency Options \hfill Week 6 \hfill \insertframenumber\end{beamercolorbox}}}
\lstset{basicstyle=\ttfamily\scriptsize,backgroundcolor=\color{codebg},
  keywordstyle=\color{ncsured},commentstyle=\color{ncsugray},
  stringstyle=\color{ncsuteal},breaklines=true,columns=fullflexible,
  showstringspaces=false,frame=single,rulecolor=\color{ncsugray}}

\title{Week 6: Options Pricing \& the Volatility Surface}
\subtitle{Greeks, no-arbitrage, and SVI/SABR calibration}
\date{North Carolina State University \,\textbar\, Summer 2026}
\titlegraphic{
\includegraphics[width=4.2cm]{/Users/dendisuhubdy/Github/ncsu/assets/ncsu_logo.png}}
\begin{document}
\begin{frame}[plain]\titlepage\begin{center}\tiny Industry Mentor: Dendi Suhubdy (Bitwyre)  |  Guest: Fenni Kang (Coincall)\end{center}\end{frame}
\begin{frame}{Contents}
\small\begin{enumerate}\setlength{\itemsep}{3pt}
  \item Learning objectives
  \item Lecture 1: Greeks \& the Surface
  \item Lecture 2: SVI \& SABR Calibration
  \item Illustration
  \item Derivation \& rationale
  \item Getting the data from Coincall
  \item Code: week6\_bs\_autograd.py
  \item Glossary of key terms
  \item Reading
  \item References
\end{enumerate}
\end{frame}

\begin{frame}{Learning Objectives}
\begin{itemize}
  \item Recall the Greeks, including vanna and volga
  \item State the no-arbitrage constraints on the surface
  \item Calibrate SVI and SABR to Coincall data
\end{itemize}
\end{frame}

\begin{frame}{Lecture 1: Greeks \& the Surface}
\begin{itemize}
  \item Black-Scholes and the Greek vector
  \item Vanna, volga, and volatility risk
  \item Calendar and butterfly arbitrage constraints
\end{itemize}
\end{frame}

\begin{frame}{Lecture 2: SVI \& SABR Calibration}
\begin{itemize}
  \item Raw and arbitrage-free SVI
  \item SVI by nonlinear least squares (autograd Jacobians)
  \item The SABR / Hagan expansion
\end{itemize}
\end{frame}

\begin{frame}{Illustration}
\begin{center}\includegraphics[width=0.86\textwidth]{/Users/dendisuhubdy/Github/ncsu/weekly_presentations/figures/week06.png}\end{center}
\end{frame}

\section*{Derivation}
\begin{frame}{Derivation \& Rationale}
\[w(k)=a+b\big(\rho(k-m)+\sqrt{(k-m)^2+\varsigma^2}\big)\]

\small\textbf{\color{ncsured}Why:} Raw SVI for total implied variance $w=\sigma^2 T$ in log-moneyness $k$. The hyperbola shape is chosen because it matches observed smiles with only 5 parameters and has closed-form slopes for the no-arbitrage checks.
\end{frame}

\begin{frame}{Derivation \& Rationale}
\[\text{butterfly: } g(k)\ge 0,\qquad \text{calendar: } \partial_T w(k,T)\ge 0\]

\small\textbf{\color{ncsured}Why:} These inequalities are exactly 'no free option-portfolio arbitrage': a negative $g(k)$ implies a negative risk-neutral density, and decreasing total variance in $T$ implies a calendar arbitrage. We enforce them so the fitted surface is tradeable.
\end{frame}

\begin{frame}{Getting the Data from Coincall}
\begin{itemize}
  \item Pull a full option chain with mark IV and Greeks per strike:
  \begin{itemize}\item \texttt{\small Chain:   GET /open/option/get/v1/\{index\}        (index = BTC/ETH underlying)}\end{itemize}
  \begin{itemize}\item \texttt{\small Detail:  GET /open/option/detail/v1/\{symbol\}    (mark, IV, delta/gamma/vega/theta/rho)}\end{itemize}
  \begin{itemize}\item \texttt{\small Symbols: GET /open/option/getInstruments/BTC}\end{itemize}
  \item Offline: coincall\_btc\_options\_chain\_snapshots\_2025Q4.parquet (5-min snapshots)
  \item A baseline coincall\_calibrated\_svi\_surfaces\_2025Q4.parquet is provided
\end{itemize}
\vspace{2pt}\tiny\textit{Cross-check exact paths at https://docs.coincall.com/}
\end{frame}

\begin{frame}[fragile]{Code --- week6\_bs\_autograd.py}
\begin{lstlisting}
import torch
def bs_call(S, K, T, r, sig):
    d1 = (torch.log(S/K) + (r + 0.5*sig**2)*T) / (sig*torch.sqrt(T))
    d2 = d1 - sig*torch.sqrt(T)
    N = torch.distributions.Normal(0., 1.).cdf
    return S*N(d1) - K*torch.exp(-r*T)*N(d2)

S,K,T,r,sig = [torch.tensor(x, requires_grad=True)
               for x in (100., 100., 1.0, 0.0, 0.6)]   # crypto IV ~ 60%
price = bs_call(S,K,T,r,sig)
vega,  = torch.autograd.grad(price, sig, create_graph=True)
vanna, = torch.autograd.grad(vega, S, create_graph=True)   # d vega / dS
volga, = torch.autograd.grad(vega, sig)                    # d vega / d sig
print(float(price), float(vega), float(vanna), float(volga))
\end{lstlisting}
\end{frame}
\begin{frame}{Glossary --- Key Terms}
\footnotesize\begin{description}
\item[\color{ncsured}Implied volatility] the volatility that matches the market option price under Black-Scholes
\item[\color{ncsured}Volatility smile] implied volatility as a function of strike / moneyness
\item[\color{ncsured}SVI] Stochastic Volatility Inspired parameterization of total variance
\item[\color{ncsured}SABR] stochastic alpha-beta-rho model with the Hagan implied-vol expansion
\item[\color{ncsured}Greeks] price sensitivities: delta, gamma, vega, theta, rho
\item[\color{ncsured}Vanna / Volga] second-order Greeks: d vega / dS and d vega / d sigma
\item[\color{ncsured}Calendar / butterfly arbitrage] no-arbitrage constraints the surface must satisfy
\end{description}
\end{frame}

\begin{frame}{Reading}
\begin{itemize}
  \item Gatheral (2006), Ch. 1-5; Gatheral-Jacquier (2014); Hagan et al. (2002)
\end{itemize}
\end{frame}

\begin{frame}{References}
\scriptsize\begin{itemize}\setlength{\itemsep}{3pt}
  \item Gatheral, J. (2006). The Volatility Surface: A Practitioner's Guide. Wiley.
  \item Gatheral, J., \& Jacquier, A. (2014). Arbitrage-free SVI volatility surfaces. Quantitative Finance, 14(1), 59-71.
  \item Hagan, P. S., Kumar, D., Lesniewski, A. S., \& Woodward, D. E. (2002). Managing smile risk. Wilmott Magazine, Sept., 84-108.
  \item Hull, J. C. (2021). Options, Futures, and Other Derivatives (11th ed.). Pearson.
\end{itemize}
\end{frame}

\end{document}